\chapter{Stochastic Calculus}\label{ch:iv:stochastic-calculus}

``What is $d(W_t^2)$?'' The candidate writes $2W_t\,dW_t$ and stops. The interviewer asks for the expectation of
$W_t^2$, and the missing term appears on its own: integrating the candidate's answer gives a quantity of zero
expectation, while $W_t^2$ is non-negative and not identically zero. Stochastic calculus questions are asked mostly of
bank quants and researchers, and they test a short list of reflexes: Itô's formula with its second-order term,
martingales and optional stopping, the reflection principle, a change of measure, and the link between an expectation
and a partial differential equation. The theory is One Quant Book~4's (chapters 2 to~5); this chapter is its use at
the whiteboard.

\section{Itô's formula as a calculation rule}

For a twice differentiable $f$ and a Brownian motion $W$, $df(t, W_t) = \big(\partial_t f + \tfrac12\partial_{xx} f\big)dt
+ \partial_x f\,dW_t$. For an Itô process $dX = \mu\,dt + \sigma\,dW$ the second-order term is $\tfrac12\sigma^2
\partial_{xx} f\,dt$. The rule to remember is that $(dW)^2 = dt$: the quadratic variation of Brownian motion is not
negligible (One Quant Book~4, chapter~3).

\begin{method}[Is it a martingale?]\label{met:iv:stochastic-calculus:martingale}
\begin{enumerate}
\item Apply Itô's formula to the process.
\item Collect the $dt$ terms: the process is a local martingale exactly when they vanish.
\item For a true martingale check integrability (for example, bounded coefficients on $[0,T]$, or Novikov's condition
for exponentials).
\item If there is a drift, subtract its integral: $W_t^2 - t$ and $W_t^3 - 3\int_0^t W_s\,ds$ are martingales.
\end{enumerate}
\end{method}

\begin{example}[The exponential martingale and the lognormal median]\label{ex:iv:stochastic-calculus:exp}
$d\exp(\sigma W_t - \sigma^2 t/2)$ has drift $-\sigma^2/2 + \sigma^2/2 = 0$: it is a martingale, and $\E[e^{\sigma W_t}] =
e^{\sigma^2 t/2}$. For $dS = \mu S\,dt + \sigma S\,dW$, $d\log S = (\mu - \sigma^2/2)\,dt + \sigma\,dW$, so the mean of $S_T$
grows at $\mu$ and its median at $\mu - \sigma^2/2$. With $\mu = 8\%$ and $\sigma = 40\%$ over ten years the mean is $2.23$
times the start and the median $1.00$: half the paths end below where they began.
\end{example}

\section{Martingales, hitting times and the reflection principle}

Exit problems for Brownian motion reduce to optional stopping on the right martingale, exactly as for the random walk
of \cref{ch:iv:probability-ii}.

\begin{proposition}[Exit from an interval]\label{prop:iv:stochastic-calculus:exit}
For $X_t = \mu t + \sigma W_t$ from 0, the probability of hitting $+a$ before $-b$ is $b/(a+b)$ if $\mu = 0$ and
$\big(1 - e^{2\mu b/\sigma^2}\big)/\big(e^{-2\mu a/\sigma^2} - e^{2\mu b/\sigma^2}\big)$ otherwise; for $\mu = 0$, $\sigma = 1$
the expected exit time is $ab$.
\end{proposition}

\begin{proof}
$W_t$ and $W_t^2 - t$ are martingales when $\mu = 0$; $\exp(-2\mu X_t/\sigma^2)$ is one otherwise. Optional stopping at the
exit time, which has finite expectation, gives the results.
\end{proof}

\begin{proposition}[Reflection principle]\label{prop:iv:stochastic-calculus:reflection}
For $a > 0$, $\P(\max_{s \le t} W_s \ge a) = 2\,\P(W_t \ge a)$.
\end{proposition}

\begin{proof}
Reflect the path after it first reaches $a$: by the strong Markov property the reflected path is again a Brownian motion,
and it ends above $a$ exactly when the original ends below; the paths that reach $a$ split evenly between ending above
and below (One Quant Book~4, chapter~2).
\end{proof}

\begin{omfigure}
\begin{tikzpicture}
\begin{axis}[width=10cm,height=5.4cm,xlabel={\small barrier $a$},ylabel={\small $\P(\max_{s\le1} W_s \ge a)$},
  xmin=0.25,xmax=2.5,ymin=0,ymax=0.85,
  tick label style={font=\footnotesize},grid=major,grid style={gray!20},table/col sep=comma,
  legend cell align=left,legend style={font=\scriptsize,at={(0.5,-0.32)},anchor=north,
  /tikz/every even column/.append style={column sep=8pt}},legend columns=3]
\addplot[omDef,thick] table[x=a,y=exact]{figdata/interviews/16-stochastic-calculus/reflect.csv};
\addplot[only marks,mark=o,mark size=1.6pt,omThm] table[x=a,y=sim100]{figdata/interviews/16-stochastic-calculus/reflect.csv};
\addplot[only marks,mark=triangle,mark size=1.8pt,omProp] table[x=a,y=sim1600]{figdata/interviews/16-stochastic-calculus/reflect.csv};
\legend{exact (reflection),100 steps,1\,600 steps}
\end{axis}
\end{tikzpicture}

\omcaption{The chance that Brownian motion on $[0,1]$ reaches the barrier $a$: the reflection principle against two
simulations of 20\,000 paths. Monitoring at discrete steps misses crossings between steps, so simulations underestimate;
dividing the step by sixteen closes most of the gap. Data: \texttt{fig\_iv\_reflect.py}.}\label{fig:iv:stochastic-calculus:reflect}
\end{omfigure}

The reflection principle gives the law of the running maximum and of the first hitting time; it also explains why a
Monte Carlo price of a barrier option monitored on a coarse grid is biased (\cref{fig:iv:stochastic-calculus:reflect}),
which is the kind of observation an interviewer rewards.

\section{Changes of measure in one line}

Girsanov's theorem (One Quant Book~4, chapter~5) says that under the measure $\mathbb Q$ with density $\exp(-\theta W_T -
\theta^2 T/2)$ the process $W_t + \theta t$ is a Brownian motion. In an interview it is used to move a drift: with $dS =
\mu S\,dt + \sigma S\,dW$ and a rate $r$, the choice $\theta = (\mu - r)/\sigma$ (the market price of risk) makes the
discounted price a martingale, which is the risk-neutral measure. Choosing the stock as numeraire instead gives the
measure under which the asset-or-nothing digital is priced.

\begin{example}[Two digitals by two measures]\label{ex:iv:stochastic-calculus:digitals}
A cash-or-nothing digital paying 1 if $S_T > K$ is worth $e^{-rT}\Phi(d_2)$ (the risk-neutral probability), and an
asset-or-nothing digital paying $S_T$ if $S_T > K$ is worth $S_0\Phi(d_1)$ (the probability under the stock measure).
A call is the difference: asset-or-nothing minus $K$ cash-or-nothing, which is the Black--Scholes formula read off in
one line (One Quant Book~5, chapter~3).
\end{example}

\section{From an SDE to a PDE and back}

The Feynman--Kac formula links expectations and PDEs: if $dX = \mu(t,X)\,dt + \sigma(t,X)\,dW$, then $u(t,x) =
\E[g(X_T) \mid X_t = x]$ solves $\partial_t u + \mu\partial_x u + \tfrac12\sigma^2\partial_{xx} u = 0$ with $u(T, x) =
g(x)$ (One Quant Book~4, chapter~4, on the generator). Interviews use it in both directions: to guess a solution of a PDE
as an expectation, and to check an expectation by plugging it into the PDE.

\begin{method}[Checking a proposed expectation]\label{met:iv:stochastic-calculus:fk}
\begin{enumerate}
\item Write the generator of the process: $\mathcal L u = \mu u_x + \tfrac12\sigma^2 u_{xx}$.
\item Check that the proposed $u$ satisfies $u_t + \mathcal L u = 0$ and the terminal condition.
\item Check a limit: at $t = T$, or for $\sigma \to 0$, the answer must reduce to the obvious one.
\end{enumerate}
\end{method}

\section{Worked answers}

\begin{example}[A mean-reverting spread]\label{ex:iv:stochastic-calculus:ou}
\emph{``A spread follows $dX_t = -\kappa X_t\,dt + \sigma\,dW_t$ with $\kappa = 0.1$ a day and $\sigma = 1$ basis point per
square-root day. What are its half-life and its long-run standard deviation?''} The mean obeys $d\E[X_t] = -\kappa\E[X_t]\,dt$,
so it decays as $e^{-\kappa t}$ and halves in $\ln 2/\kappa \approx 6.9$ days. For the variance, Itô's formula on $X^2$
gives $d\E[X_t^2] = (-2\kappa\E[X_t^2] + \sigma^2)\,dt$, whose stationary point is $\sigma^2/(2\kappa) = 5$, a standard
deviation of $\sqrt5 \approx 2.24$ basis points. The two numbers are what a pairs trader sizes and times entries with: a
spread two standard deviations out, 4.5 basis points, is expected to be halfway back in a week. \emph{Check:} as $\kappa \to
0$ the stationary variance blows up, as it must for a random walk.
\end{example}

\begin{example}[How long a driftless P\&L stays under water]\label{ex:iv:stochastic-calculus:arcsine}
\emph{``A strategy's cumulative P\&L is a driftless Brownian motion started at zero. What is the chance it spends at least
90\% of the year below zero?''} By Lévy's arcsine law, the fraction of $[0, 1]$ spent positive has distribution function
$\tfrac2\pi\arcsin\sqrt{x}$, so the chance of at most 10\% of the time positive is $\tfrac2\pi\arcsin\sqrt{0.1} \approx
0.20$. One in five worthless strategies, or skilful strategies with no edge this year, spend nine-tenths of the year under
water, and as many spend nine-tenths above it. The density piles up at the ends, not in the middle: long spells on one
side are the rule for a random walk, which is why a run of good months says less than it feels.
\end{example}

\section{Question bank}

\begin{interviewq}[$\star$ \iqroles{bank, researcher} \iqfirm{bank}]\label{iq:iv:stochastic-calculus:1}
Compute $d(W_t^2)$ and use it to find $\E[W_t^2]$.
\end{interviewq}

\begin{interviewq}[$\star$ \iqroles{bank, researcher} \iqfirm{bank}]\label{iq:iv:stochastic-calculus:2}
Is $W_t^3$ a martingale? If not, what must be subtracted to make it one?
\end{interviewq}

\begin{interviewq}[$\star$ \iqroles{researcher, trader} \iqfirm{systematic fund}]\label{iq:iv:stochastic-calculus:3}
What is the probability that a Brownian motion is positive at both $t = 1$ and $t = 2$?
\end{interviewq}

\begin{interviewq}[$\star$ \iqroles{bank} \iqfirm{bank}]\label{iq:iv:stochastic-calculus:4}
What is the distribution of $\int_0^T W_s\,ds$?
\end{interviewq}

\begin{interviewq}[$\star\star$ \iqroles{bank, researcher} \iqfirm{market maker}]\label{iq:iv:stochastic-calculus:5}
Starting from 0, what is the chance that a Brownian motion hits $+1$ before $-2$, and how long does it take on average to hit
either?
\end{interviewq}

\begin{interviewq}[$\star\star$ \iqroles{researcher, trader} \iqfirm{systematic fund}]\label{iq:iv:stochastic-calculus:6}
$X_t = 0.5t + W_t$. What is the chance it hits $+1$ before $-1$?
\end{interviewq}

\begin{interviewq}[$\star\star$ \iqroles{bank, risk} \iqfirm{bank}]\label{iq:iv:stochastic-calculus:7}
What is the probability that a standard Brownian motion exceeds 1 at some time in $[0, 1]$? A Monte Carlo with 50 steps
gives a smaller number; why, and how do you fix it without more steps?
\end{interviewq}

\begin{interviewq}[$\star\star$ \iqroles{trader, risk} \iqfirm{asset manager}]\label{iq:iv:stochastic-calculus:8}
A stock follows geometric Brownian motion with drift 8\% and volatility 40\%. Over ten years, find the expected value and the
median of its price relative to today. What do you tell an investor who hears only the first number?
\end{interviewq}

\begin{interviewq}[$\star\star$ \iqroles{bank} \iqfirm{bank}]\label{iq:iv:stochastic-calculus:9}
A stock has drift 8\% and volatility 20\% and the rate is 3\%. What change of measure makes the discounted stock a
martingale, and what is the drift of $W$ under it?
\end{interviewq}

\begin{interviewq}[$\star\star\star$ \iqroles{bank, trader} \iqfirm{bank}]\label{iq:iv:stochastic-calculus:10}
With $S_0 = K = 100$, $\sigma = 20\%$, $T = 1$ and zero rates, price the cash-or-nothing digital paying 1 and the
asset-or-nothing digital paying $S_T$ if $S_T > K$. Deduce the call price.
\end{interviewq}

\begin{interviewq}[$\star\star\star$ \iqroles{bank, researcher} \iqfirm{bank}]\label{iq:iv:stochastic-calculus:11}
Solve $\partial_t u + \tfrac12\partial_{xx}u = 0$ with $u(T, x) = x^2$ by writing $u$ as an expectation, and check it.
\end{interviewq}

\begin{interviewq}[$\star\star\star$ \iqroles{researcher, bank} \iqfirm{systematic fund}]\label{iq:iv:stochastic-calculus:12}
For $X_t = \mu t + W_t$ with $\mu > 0$, what is the expected time to reach $+1$? What happens as $\mu \to 0$, and is that
consistent with Brownian motion reaching $+1$ with probability one?
\end{interviewq}

\begin{interviewq}[$\star\star\star$ \iqroles{researcher, bank} \iqfirm{any}]\label{iq:iv:stochastic-calculus:13}
Let $\tau$ be the first time $W$ hits 1. Then $W_\tau = 1$, yet $W$ is a martingale started at 0. Why does optional
stopping fail? Add a lower barrier at $-b$ and recover a correct statement.
\end{interviewq}

\begin{omsources}
\item One Quant Book~4, chapters 2--5 (Brownian motion, Itô calculus, SDEs, Girsanov); One Quant Book~5, chapter~3
(Black--Scholes).
\item I.~Karatzas and S.~E.~Shreve, \emph{Brownian Motion and Stochastic Calculus}, Springer, 1988.
\end{omsources}
